\begin{tikzpicture}[font=\scriptsize, >=latex]
% book map
\foreach \x/\t/\d in {
  0.2/{Part I (ch1--4)}/{source $\to$ machine code},
  3.4/{Part II (ch5--8)}/{the RV32I capstone},
  6.6/{Part III (ch9--11)}/{memory hierarchy}}{
  \draw[thick, fill=accentsoft] (\x,3.4) rectangle (\x+2.9,4.8);
  \node[align=center, font=\footnotesize] at (\x+1.45,4.4) {\t};
  \node[align=center, font=\tiny] at (\x+1.45,3.9) {\d};
}
\foreach \x/\t/\d in {
  1.8/{Part IV (ch12--14)}/{the OS: processes, files},
  5.0/{Part V (ch15--17)}/{performance craft},
  8.2/{Part VI (ch18)}/{runtimes}}{
  \draw[thick, fill=graybg] (\x,1.6) rectangle (\x+2.9,3.0);
  \node[align=center, font=\footnotesize] at (\x+1.45,2.6) {\t};
  \node[align=center, font=\tiny] at (\x+1.45,2.1) {\d};
}
\draw[->, thick] (3.1,4.1) -- (3.4,4.1);
\draw[->, thick] (6.3,4.1) -- (6.6,4.1);
\draw[->, thick] (7.9,3.4) -- (6.8,3.0);
\draw[->, thick] (4.7,3.0) -- (4.0,3.4);
\draw[->, thick] (6.5,3.0) -- (6.6,3.4);
\node[align=left, text=slate, text width=10cm] at (5.2,0.6)
{The path this chapter starts: build the real machine around the
toy CPU you already made, then measure what the folklore got
wrong.};
\end{tikzpicture}
